\section{Safe envelopes for signed accelerated candidates}
\label{sec:safe-envelope}

Acceleration may overshoot, so its candidate cannot itself control support.
The next lemma converts any signed restricted candidate into a certified lower
point.

For $U\subseteq S_\rho$, let $x_U^\star$ solve the restricted nonnegative
quadratic program
\begin{equation}
  x_U^\star\in\arg\min_{z_U\geq0}
  \left\{\frac12z_U^\top Q_{UU}z_U-c_{\rho,U}^\top z_U\right\},
  \qquad x_{U^c}^\star=0.
  \label{eq:restricted-QP}
\end{equation}
The Stieltjes obstacle comparison principle gives
\begin{equation}
  0\leq x_U^\star\leq x_\rho^\star.
  \label{eq:restricted-comparison}
\end{equation}
Indeed, $x_{\rho,U}^\star$ is a supersolution of the restricted KKT system
because $U\subseteq S_\rho$ and
$Q_{UU}x_{\rho,U}^\star-c_{\rho,U}
=-Q_{UU^c}x_{\rho,U^c}^\star\geq0$.

\begin{theorem}[Signed candidate to safe lower envelope]
\label{thm:safe-envelope}
Let $z_U\in\R^U$ be arbitrary and define
\begin{equation}
 q_U:=c_{\rho,U}-Q_{UU}z_U,
 \qquad
 \delta_U:=\frac1\alpha
  \norm{D_U^{-1/2}(q_U)_-}_\infty.
 \label{eq:envelope-delta}
\end{equation}
Set
\begin{equation}
 \ell_U:=\left[z_U-\delta_UD_U^{1/2}\one\right]_+,
 \qquad \ell_{U^c}:=0.
 \label{eq:safe-envelope}
\end{equation}
Then
\begin{equation}
  0\leq\ell\leq x_U^\star\leq x_\rho^\star.
  \label{eq:envelope-order}
\end{equation}
\end{theorem}

\begin{proof}
Let $q_U^\star:=c_{\rho,U}-Q_{UU}x_U^\star\leq0$ be the restricted KKT
residual.  Since $Q_{UU}^{-1}\geq0$,
\begin{align*}
 x_U^\star-z_U
 &=Q_{UU}^{-1}(q_U-q_U^\star)
 \geq-Q_{UU}^{-1}(q_U)_-.
\end{align*}
The full eigenvector identity in \cref{eq:Q-properties} and the signs of the
cut entries imply
\begin{equation}
 Q_{UU}D_U^{1/2}\one
 =\alpha D_U^{1/2}\one-Q_{UU^c}D_{U^c}^{1/2}\one
 \geq\alpha D_U^{1/2}\one.
\end{equation}
Thus
$Q_{UU}^{-1}D_U^{1/2}\one\leq\alpha^{-1}D_U^{1/2}\one$.
By \cref{eq:envelope-delta},
$Q_{UU}^{-1}(q_U)_-\leq\delta_UD_U^{1/2}\one$, and hence
$z_U-\delta_UD_U^{1/2}\one\leq x_U^\star$.  Taking the positive part preserves
the inequality because $x_U^\star\geq0$.  The last comparison is
\cref{eq:restricted-comparison}.
\end{proof}

The accelerator is therefore free to use signed momentum in $z$; only
$\ell$ is allowed to govern expansion.
